\subsection{科恩环}

定义 2．——设 A 是一个分离且完备的局部环，其剩余域的特征为 p。称 A 的一个子环 C 为 A 的一个科恩子环，如果 C 是 p-环且满足 \(\mathbf{A} = \mathfrak{m}_{\mathbf{A}} + \mathbf{C}\)（即 \(\mathbf{A}/\mathfrak{m}_{\mathbf{A}} = \mathbf{C}/(\mathfrak{m}_{\mathbf{A}} \cap \mathbf{C})\)）。

若 C 是 A 的一个科恩子环，则 C 的理想 \(m_{A} \cap C\) 是极大理想，因而等于 pC。所以从 \(\kappa_{C} = C/pC\) 到 \(\kappa_{A} = A/m_{A}\) 的典范映射是域同构。

例．——设 C 是 p-环。形式幂级数环 \(A = C[[T_{1}, \ldots, T_{n}]]\) 是一个诺特局部环，分离且完备，其极大理想由序列 (p, T₁, \ldots, Tₙ) 生成。显然 C 是 A 的一个科恩子环。特别地，当 C 等于 Zₚ、Z/pⁿZ 或特征为 p 的域时也适用。

定理 1．— 设 A 是一个局部环，分离且完备，其剩余域 k 的特征为 p。设 π 是从 A 到 k 的典范映射，并设 S 是 A 的一个子集，使 π 将 S 双射到 k 的一个 p-基上（A，V，p. 95）。

\begin{enumerate}
\def\labelenumi{\alph{enumi})}
\item
  存在唯一一个包含 S 的 A 的科恩子环 C。
\item
  A 的子环 C 是闭的，且 C 上的 \(p\)C-进拓扑由 A 的 \(m_{A}\)-进拓扑诱导。
\item
  A 的每个包含 S 且满足 \(A = A' + m_{A}\) 的闭子环 A' 都包含 C。
\end{enumerate}

\begin{enumerate}
\def\labelenumi{\Alph{enumi})}
\tightlist
\item
  特例：\(\mathfrak{m}_{\mathbf{A}}\) 是幂零的
\end{enumerate}

令 \(n\) 为满足 \(\mathfrak{m}_{\mathbf{A}}^{n+1}=\{0\}\) 的正整数。若 \(\Phi_{n}\) 是第 \(n\) 个 Witt 多项式（§ 1，第 1 号），则映射 \(u:[a_{0},..., a_{n}]\mapsto\Phi_{n}(a_{0},..., a_{n})\) 是从 \(\mathbf{W}_{n+1}(\mathbf{A})\) 到 \(\mathbf{A}\) 的环同态（§ 1，第 7 号）。令 \(\mathbf{B}_{n}\) 为 \(u\) 的像，并令 \(\mathbf{C}_{n}\) 为 \(\mathbf{A}\) 的由 \(\mathbf{B}_{n}\cup\mathbf{S}\) 生成的子环。

引理 1．——设 A' 是包含 S 的 A 的一个子环。A' 包含 \(\mathbf{C}_{n}\) 的充要条件是 \(A' + m_{A} = A\)。

有 \(pA \subset m_A\)，且 \(B_n\) 由形如 \(a_0^{p^n} + pa_1^{p^{n-1}} + \cdots + p^n a_n\) 的元素组成，其中 \(a_0, \ldots, a_n\) 属于 A。因此 \(\pi(B_n) = k^{p^n}\)，从而 \(\pi(C_n) = k^{p^n}[\pi(S)]\)。但由于 \(\pi(S)\) 是一个 \(p\)-基（属于 \(k\)），有 \(k = k^{p^n}[\pi(S)]\)（A，V，p. 96），从而 \(\pi(C_n) = k\)，即 \(C_n + m_A = A\)。

设 A' 是包含 S 的 A 的一个子环。若 A' 包含 \(C_{n}\)，则有

\[
\mathrm{A} ^ {\prime} + \mathrm{m} _ {\mathrm{A}} \supset \mathrm{C} _ {n} + \mathrm{m} _ {\mathrm{A}} = \mathrm{A}, \quad \mathrm{whence} \quad \mathrm{A} ^ {\prime} + \mathrm{m} _ {\mathrm{A}} = \mathrm{A}.
\]

反之，设 \(\mathbf{A}' + \mathfrak{m}_{\mathbf{A}} = \mathbf{A}\)。令 \(a_0, \ldots, a_n\) 为 \(\mathbf{A}\) 中的元素；根据假设，存在元素 \(a_0', \ldots, a_n'\) 属于 \(\mathbf{A}'\)，使得 \(a_i \equiv a_i'\) 模 \(\mathfrak{m}_{\mathbf{A}}\)（对 \(0 \leqslant i \leqslant n\)）。根据 § 1 第 1 号的命题 1 及假设 \(\mathfrak{m}_{\mathbf{A}}^{n+1} = \{0\}\)，有 \(\Phi_n(a_0, \ldots, a_n) = \Phi_n(a_0', \ldots, a_n') \in \mathbf{A}'\)，从而 \(\mathbf{B}_n \subset \mathbf{A}'\)。由于 \(\mathbf{C}_n\) 是由 \(\mathbf{B}_n \cup \mathbf{S}\) 生成的环，有 \(\mathbf{C}_n \subset \mathbf{A}'\)。

在所有包含 S 且满足 \(A' + m_{A} = A\) 的 A 的子环 A' 所成的集合中，根据引理 1 存在一个最小元素 C，并且有 \(C_{n} = C\)，对于每个整数 \(n \geqslant 0\) 满足 \(m_{A}^{n+1} = \{0\}\) 时均如此。

按构造有 \(C + m_A = A\)，且 \(p1_C\) 是幂零的。显然 \(pC \subset C \cap m_A\)，因此下面的引理 2 表明 \(pC\) 是 \(C\) 的极大理想，从而 \(C\) 是 \(A\) 的一个科恩子环。

引理 2．——有 \(C \cap m_{A} \subset pC\)。

取整数 \(m \geqslant 1\)，使得 \(\mathfrak{m}_{\mathbf{A}}^{m} = \{0\}\)，从而 \(\mathbf{C} = \mathbf{C}_{m} = \mathbf{C}_{m-1}\)。令 \(\Lambda\) 为 \(\mathbf{N}^{(\mathbf{S})}\) 的一部分，它由有限支撑的整数族 \((\alpha_{s})_{s \in \mathbf{S}}\) 组成，并满足 \(0 \leqslant \alpha_{s} < p^{m}\)，对每个 \(s \in \mathbf{S}\) 均如此。由于 \(\mathbf{B}_{m}\) 包含 \(s^{p^{m}} = \Phi_{m}(s, 0, ..., 0)\)（对每个 \(s \in \mathbf{S}\)），单项式 \(Z_{\alpha} = \prod_{s \in \mathbf{S}} s^{\alpha_{s}}\)（其中 \(\alpha\) 遍历 \(\Lambda\)）生成 \(\mathbf{C}_{m}\)，作为 \(\mathbf{B}_{m}\)-模。此外，根据公式

\[
\Phi_ {m} (a _ {0}, \dots , a _ {m}) = a _ {0} ^ {p ^ {m}} + p \Phi_ {m - 1} (a _ {1}, \dots , a _ {m}),
\]

 \(B_{m}\) 的每个元素都形如 \(a^{p^{m}} + pb\)，其中 \(a \in A\) 且 \(b \in B_{m-1}\)。因此 \(C = C_{m}\) 的每个元素都形如

\[
x = \sum_ {\alpha \in \Lambda} c _ {\alpha} ^ {p ^ {m}} Z _ {\alpha} + p y\tag{1}
\]

其中 \(c_{\alpha} \in \mathbf{A}\) 对每个 \(\alpha \in \Lambda\) 均成立，且 \(y \in \mathbf{C}_{m-1} = \mathbf{C}\)。若 \(x\) 属于 \(\mathbf{C} \cap \mathfrak{m}_{\mathbf{A}}\)，则有 \(\pi(x) = 0\)，从而 \(\sum_{\alpha \in \Lambda} \pi(c_{\alpha})^{p^m} \pi(Z_{\alpha}) = 0\)。由于 \(\pi(\mathbf{S})\) 是一个 \(p\)-基（属于 \(k\)），根据 A，V，p. 96，有 \(\pi(c_{\alpha}) = 0\) 对每个 \(\alpha \in \Lambda\) 均成立。于是 \(c_{\alpha} \in \mathfrak{m}_{\mathbf{A}}\)，从而 \(c_{\alpha}^m = 0\)，当然也有 \(c_{\alpha}^{p^m} = 0\)。由 (1) 得 \(x = py\)，从而得到引理 2。

有 \(p^m\mathbf{C} = \mathfrak{m}_{\mathbf{A}}^m = \{0\}\)（当 \(m\) 足够大时），所以断言 \(b)\) 显然成立。断言 \(c)\) 由引理 1 得出。若 \(\mathbf{C}'\) 是包含 S 的 A 的一个科恩子环，则根据引理 1 有 \(\mathbf{C}' \supset \mathbf{C}\)。但由于 C 到 \(\mathbf{C}'\) 的包含映射诱导出从 \(\kappa_{\mathbf{C}}\) 到 \(\kappa_{\mathbf{C}'}\) 的同构，有 \(\mathbf{C} = \mathbf{C}'\)（第 1 号，命题 2 b)），这就完成了 \(a\) 的证明。

\noindent\textbf{B) 一般情形.}\par

对于每个整数 \(n \geqslant 0\)，记 \(\mathbf{A}_n\) 为局部环 \(\mathbf{A}/\mathfrak{m}_{\mathbf{A}}^{n+1}\)，\(\mathfrak{m}_n = \mathfrak{m}_{\mathbf{A}}/\mathfrak{m}_{\mathbf{A}}^{n+1}\) 为其极大理想，并记 \(\pi_n\) 为从 A 到 \(\mathbf{A}_n\) 的典范同态。根据 A)，存在唯一的科恩子环 \(\mathbf{C}_n\)，它是 \(\mathbf{A}_n\) 的且包含 \(\pi_n(\mathbf{S})\)。当 \(0 \leqslant n \leqslant m\) 时，记 \(\pi_{n,m}\) 为从 \(\mathbf{A}_m\) 到 \(\mathbf{A}_n\) 的典范同态。根据第 1 号的命题 1 的推论 2，\(\pi_{n,m}(\mathbf{C}_m)\) 是一个 \(p\)-环；有 \(\pi_{n,m}(\mathbf{C}_m) + \mathfrak{m}_n = \mathbf{A}_n\)，所以 \(\pi_{n,m}(\mathbf{C}_m)\) 等于 \(\mathbf{C}_n\)，即 \(\mathbf{A}_n\) 的科恩子环。根据第 1 号的命题 3，子环 \(\varprojlim \mathbf{C}_n\) 位于 \(\varprojlim \mathbf{A}_n\) 中，是一个 \(p\)-环。置 \(\mathbf{C} = \bigcap_{n \in \mathbb{N}} \pi_n^{-1}(\mathbf{C}_n)\)。由于 C 是 \(\varprojlim \mathbf{C}_n\) 在同构 \(a \mapsto (\pi_n(a))_{n \in \mathbb{N}}\) 下的逆像，而该同构将 A 映到 \(\varprojlim \mathbf{A}_n\)，故 C 是 A 的闭子环，也是一个 \(p\)-环。有 \(\pi_n(\mathbf{C}) = \mathbf{C}_n\)（对每个 \(n \in \mathbb{N}\)），特别地 \(\pi_0(\mathbf{C}) = \mathbf{A}_0\)，即 \(\pi(\mathbf{C}) = k\)。所以 C 是 A 的一个科恩子环。

对于每个整数 \(n \geqslant 0\)，置 \(J_{n} = C \cap m_{A}^{n}\)。由于局部环 A 是分离的，有 \(\bigcap_{n \in \mathbb{N}} J_{n} = \{0\}\)，并且根据 \(p\)-环理想的结构（第 1 号，命题 1），C 的每个形如 \(p^{k}C\) 的理想都包含某个 \(J_{n}\)。反之，\(J_{n}\) 包含 \(p^{n}C\)。所以 C 的 \(pC\)-进拓扑由 A 的 \(m_{A}\)-进拓扑诱导。这证明了 \(b\)。

设 A' 是 A 的一个闭子环，包含 S 且满足 \(A' + m_{A} = A\)。由于 A' 是闭的，有 A' = \(\bigcap_{n\in\mathbb{N}}\pi_n^{-1}(\pi_n(A'))\)。有 \(\pi_n(A') \supset \pi_n(S)\) 且 \(\pi_n(A') + m_n = A_n\)，从而根据 A) 中已经得到的结论，\(\pi_n(A') \supset C_n\)。最后，有 \(\pi_n^{-1}(\pi_n(A')) \supset \pi_n^{-1}(C_n)\)，从而 A' ⊃ C。这证明了 c)。由此得到 A) 中所述科恩子环的唯一性。

注记．——若 \(p1_{A}\) 不是幂零的（特别地，当 A 是整环且其分式域的特征为 0 时会发生这种情况），则 C 是离散赋值环，其分式域的特征为 0。

